%~Mouliné par MaN_auto v.0.21.8 2020-03-13 13:23:51
\documentclass[ALCO,ThmDefs,Unicode]{cedram}
\OneNumberAllTheorems

\TralicsDefs{\addattributestodocument{type}{corrigendum}}
\Relation{ALCO_2019__2_6_1197_0}

%\usepackage{bbm}
%\usepackage[linesnumbered,lined,ruled]{algorithm2e}
%\usepackage{mathrsfs}
%\usepackage{float}
%\usepackage{stmaryrd}

\newtheorem{innercustomthm}{Theorem}
\newenvironment{customthm}[1]
  {\renewcommand\theinnercustomthm{#1}\innercustomthm}
  {\endinnercustomthm}

\graphicspath{{./figures/}}

\newcommand*{\mk}{\mkern -1mu}
\newcommand*{\Mk}{\mkern -2mu}
\newcommand*{\mK}{\mkern 1mu}
\newcommand*{\MK}{\mkern 2mu}

\hypersetup{urlcolor=purple, linkcolor=blue, citecolor=red}

\newcommand*{\romanenumi}{\renewcommand*{\theenumi}{\roman{enumi}}}
\newcommand*{\Romanenumi}{\renewcommand*{\theenumi}{\Roman{enumi}}}
\newcommand*{\alphenumi}{\renewcommand*{\theenumi}{\alph{enumi}}}
\newcommand*{\Alphenumi}{\renewcommand*{\theenumi}{\Alph{enumi}}}


\title{Corrigendum to ``Graphs of gonality three''}

\author[\initial{I.} Aidun]{\firstname{Ivan} \lastname{Aidun}}
\address{Oberlin College\\
Department of Mathematics\\
10 N. Professor St.\\
Oberlin\\
OH 44074, USA}
\email{iaidun@oberlin.edu}

\author[\initial{F.} Dean]{\firstname{Frances} \lastname{Dean}}
\address{Williams College\\
Department of mathematics and statistics\\
33 Stetson Ct.\\
Williamstown\\
MA 01267, USA}
\email{fed1@williams.edu}

\author[\initial{R.} Morrison]{\firstname{Ralph} \lastname{Morrison}}
\address{Williams College\\
Department of mathematics and statistics\\
33 Stetson Ct.\\
Williamstown\\
MA 01267, USA}
\email{10rem@williams.edu}
\thanks{The authors were supported by NSF grants DMS1659037 and DMS1347804, and by the Williams College SMALL REU}

\author[\initial{T.} Yu]{\firstname{Teresa} \lastname{Yu}}
\address{Williams College\\
Department of mathematics and statistics\\
33 Stetson Ct.\\
Williamstown\\
MA 01267, USA}
\email{twy1@williams.edu}

\author[\initial{J.} Yuan]{\firstname{Julie} \lastname{Yuan}}
\address{University of Minnesota-Twin Cities\\
School of mathematics\\
206 Church St SE\\
Minneapolis\\
MN 55455, USA}
\email{yuanx254@umn.edu}

\keywords{graph gonality, chip-firing, tropical geometry}
\subjclass{14T05, 05C05, 05C57}

\begin{abstract}{}
Our main theorem in~\cite{original} contained a mistake regarding the equivalence of two conditions on a graph, which we correct here. Fortunately our main result is not impacted with an additional assumption called the zero-three condition.
\end{abstract}

\begin{document}
\maketitle

The following is a corrected version of the statement of our main theorem, Theorem~1.2.
We say a graph of gonality $3$ \emph{satisfies the zero-three condition} if there exists a divisor $D$ such that $\deg(D)=3$, $r(D)=1$, and for any three distinct vertices $a,b,c$ with $D\sim (a)+(b)+(c)$, we have that $a,b,c$ either share $0$ edges or share $3$ edges.

\renewcommand*{\thecdrthm}{1.2}
\begin{theo}\label{thm:maintheorem}
If $G$ is a $3$-edge-connected combinatorial graph, then the following are equivalent:
\begin{enumerate}
\item \label{1.2.1} $G$ has \textup{(}divisorial\textup{)} gonality $3$.
\item \label{1.2.2} There exists a non-degenerate harmonic morphism $\varphi : G \to T$, where $\deg(\varphi) = 3$ and $T$ is a tree.
\end{enumerate}
Moreover, if $G$ is simple and $3$-vertex-connected, and also satisfies the zero-three condition, these statements imply the following condition:
\begin{enumerate}\setcounter{enumi}{2}
\item \label{1.2.3} There exists a cyclic automorphism $\sigma: G \to G$ of order $3$ that does not fix any edge of $G$ satisfying the property that $G/\sigma$ is a tree.
\end{enumerate}
Condition~\eqref{1.2.3} implies conditions~\eqref{1.2.1} and~\eqref{1.2.2}, whether or not $G$ satisfies the zero-three condition.
\end{theo}

Similarly, we have the following corrected version of Theorem~4.1.

\renewcommand*{\thecdrthm}{4.1}
\begin{theo}\label{thm:simpletheorem}
If $G$ is a simple, $3$-vertex-connected combinatorial graph satisfying the zero-three condition, then the following are equivalent:
\begin{enumerate}
\item \label{4.1.1} $G$ has gonality $3$.
\item \label{4.1.2} There exists a non-degenerate harmonic morphism $\varphi : G \to T$, where $\deg(\varphi) = 3$ and $T$ is a tree.
\item \label{4.1.3} There exists a cyclic automorphism $\sigma: G \to G$ of order $3$ that does not fix any edge of $G$, such that $G/\sigma$ is a tree.
\end{enumerate}
The equivalence of~\eqref{4.1.1} and~\eqref{4.1.2} and the implication~\eqref{4.1.3} implies~\eqref{4.1.1} and~\eqref{4.1.2} still hold without the zero-three condition assumption.
\end{theo}

\begin{figure}[hbt]
\centering
\includegraphics[scale=1]{not_zero_three.pdf}
\caption{The wheel graph $W_5$}\label{fig:wheel}
\end{figure}

\goodbreak
To see that the zero-three condition is necessary, consider the graph in Figure~\ref{fig:wheel}, which is the wheel graph $W_5$ on $5$ vertices. It has gonality $3$: the divisor $(w_1)+(h)+(w_3)$ has positive rank; and since the graph has $K_4$ as a minor, the treewidth of the graph, and thus its gonality, is at least $3$. It is also $3$-vertex-connected. However, we claim that it does not satisfy the zero-three condition. Let $D$ be any effective rank $1$ divisor of degree~$3$ on $W_5$. By Lemma~4.2, $D$ must have support size $1$ or $3$, and it is equivalent to some divisor with support size $3$. If $W_5$ satisfies the zero-three condition, then since any three vertices have at least one edge in common, there must be a rank $1$ divisor $(a)+(b)+(c)$ on $W_5$ where $a,b,c$ are distinct and form a $K_3$ in the graph. It follows that $(w_i)+(w_{i+1})+(h)$ has rank $1$ for some $i$, where addition is done modulo $4$. However, this divisor does not have rank $1$: starting Dhar's burning algorithm from $w_{i+2}$ burns the whole graph. Thus $W_5$ must not satisfy the zero-three condition. Finally, the automorphism group of $W_5$ is the same as the automorphism group of the cycle $C_4$, namely the dihedral group of order $8$. This group does not have any elements of order $3$. Thus, $W_5$ is an example of a simple, $3$-vertex-connected graph of gonality $3$, not satisfying the zero-three condition, which does not have any automorphisms of order~$3$.


The gap in the original argument, which did not assume the zero-three condition, came when proving the map $\sigma$ was an automorphism in Proposition~4.3. It was true that $\sigma$ was well-defined and preserved adjacency relations between different classes of points under $\sim_D$, but it was omitted to prove that we could send the three points in the same class under $\sim_D$ to one another in a nontrivial way; this can be done precisely when there are zero or three edges between them. Thus, the proof needs to assume the zero-three condition, and we would add the following to the beginning of its proof: ``First consider the action of $\sigma$ on a triple $v_1,v_2,v_3$ where $D\sim(v_1)+(v_2)+(v_3)$ and $v_1,v_2,v_3$ are all distinct. We are assuming our graph satisfies the zero-three condition, so the map $\sigma$ mapping $v_1$ to $v_2$ to $v_3$ to $v_1$ preserves the connectivity of $v_1,v_2,$ and $v_3$, since either all share edges or none share edges.''

The only one of our other results or examples relying on the incorrect statements of Theorems~1.2 and 4.1 was our consideration of the Frucht graph, in Figure~8. It is no longer obvious that it cannot have gonality $3$: we only know that if it does have gonality $3$, then it does not satisfy the zero-three property. However, we have still computationally verified that the Frucht graph has gonality $4$, as claimed.

\bibliographystyle{amsplain-ac}
\nocite{*}
\bibliography{ALCO_Aidun_417}
\end{document}
